\chapter{Aksioma Bilangan Real dan Supremum}

Sebuah program mencari panjang diagonal persegi bersisi satu. Ia memelihara dua
bilangan, $a=1$ dan $b=2$, dengan $a^2-2<0<b^2-2$. Setiap langkah ia menghitung titik
tengah $m=(a+b)/2$, menilai $m^2-2$, lalu mengganti $a$ atau $b$ dengan $m$ supaya
tanda kedua ujung tetap berbeda. Setelah 40 langkah lebar selangnya di bawah $10^{-12}$,
setelah 100 langkah di bawah $10^{-30}$. Tetapi program itu tidak pernah mengeluarkan
jawaban: tidak ada langkah yang menghasilkan $x^2-2=0$ persis, dan kedua ujung selang
terus mengapit sesuatu yang bukan bilangan rasional.

Kegagalan itu bukan kesalahan algoritme. Ia menunjukkan bahwa bilangan rasional punya
lubang, dan bahwa himpunan $\{q\in\mathbb{Q}: q^2<2\}$ tidak mempunyai batas atas
terkecil di dalam $\mathbb{Q}$.

Bab 1 mengajari kita menulis bukti. Bab 2 memberi bahasa himpunan dan fungsi, sekaligus
memakai $|x|$ dan $\sup$ sebagai alat hitung tanpa mempertanyakan asalnya. Yang belum
ada adalah objek yang menjamin alat itu sah. Bab ini membangunnya, bukan dari desimal,
melainkan dari tiga kelompok pernyataan yang diterima tanpa bukti: hukum lapangan,
hukum urutan, dan kelengkapan. Sisanya, termasuk keberadaan $\sqrt{2}$, diturunkan
dari situ.

\begin{untukanak}
Bayangkan sebuah penggaris yang hanya punya tanda pecahan: setengah, seperempat,
seperdelapan, dan terus makin rapat. Tidak ada angka lain di penggaris itu. Sekarang
ukur diagonal sebuah persegi yang sisinya satu jengkal. Ujung diagonal jatuh di antara
dua tanda, dan selalu begitu, tidak peduli seberapa rapat tandanya. Kita bisa mendekat
sejauh yang kita mau, tetapi tidak pernah tepat.

Pertanyaannya: apakah panjang itu benar-benar ada? Penggarisnya jelas tidak punya
jawaban. Namun diagonal itu, kalau dijinjing atau dipasang di lantai, punya panjang
tertentu. Panjang itu tidak hilang hanya karena alat ukur kita tidak memuatnya.

Cerita yang sama terjadi pada bilangan. Kalau kita hanya punya bilangan pecahan, ada
titik di garis yang tidak punya nama. Matematikawan menutup lubang itu dengan satu
janji: setiap kumpulan bilangan yang terbatas ke atas punya atap paling rendah. Di bab
ini kita akan tahu mengapa janji itu perlu, apa akibatnya, dan bagaimana ia memberi
nama pada panjang diagonal tadi, yaitu akar dua.
\end{untukanak}

\section{Aksioma lapangan}
\tingkat{dasar}

Kita tidak mendefinisikan bilangan real sebagai desimal tak berujung: cara itu membuat
penjumlahan rumit, dan membuat definisi batas atas terkecil berputar-putar, sebab
``supremum semua pemotongan desimal'' sudah memakai bilangan yang sedang dibangun.
Lebih bersih menyebut sifatnya lebih dulu, lalu menurunkan sisanya.

Lapisan pertama bersifat aljabar. Kita mengandaikan dua operasi yang menaati aturan
yang kita pakai sejak sekolah; definisi berikut hanya memberi nama aturan itu.

\begin{definisi}[Lapangan]\index{lapangan}
Himpunan $F$ disebut \emph{lapangan} jika pada anggotanya terdefinisi penjumlahan dan
perkalian sehingga: hasil kedua operasi selalu berada di $F$; kedua operasi komutatif
dan asosiatif; ada unsur $0$ dengan $0+x=x$ dan unsur $1\neq 0$ dengan $1\cdot x=x$;
setiap $x$ punya lawan $-x$ dengan $x+(-x)=0$; setiap $x\neq 0$ punya kebalikan $1/x$
dengan $x\cdot(1/x)=1$; dan berlaku $x(y+z)=xy+xz$.
\end{definisi}

Aksioma itu tidak menyebut urutan, tidak menyebut jarak, dan tidak menyebut
kelengkapan. Karena itu himpunan lain juga memenuhinya. Bilangan rasional dengan
penjumlahan dan perkalian biasa adalah lapangan. Pertanyaan yang wajar: apa saja yang
sudah terkandung di dalam aturan itu?

Tiga pernyataan aljabar yang biasa kita hafal ternyata akibat, bukan aksioma; yang
paling sering dipakai adalah hasil kali dengan nol yang memusnahkan.

\begin{teorema}
Dalam sebarang lapangan, untuk semua $x,y$ berlaku: (a) $0\cdot x=0$;
(b) $(-x)y=-(xy)$; (c) $(-x)(-y)=xy$; (d) jika $xy=0$ maka $x=0$ atau $y=0$.
\end{teorema}

\begin{proof}
(a) Dari $0+0=0$ dan hukum distributif, $0\cdot x=(0+0)x=0\cdot x+0\cdot x$.
Tambahkan $-(0\cdot x)$ pada kedua ruas. Ruas kiri menjadi $0$, ruas kanan menjadi
$0\cdot x$. Jadi $0=0\cdot x$.

(b) Hitung $(-x)y+xy=(-x+x)y=0\cdot y=0$, memakai (a). Jadi $(-x)y$ lawan dari $xy$,
dan lawan itu tunggal.

(c) Terapkan (b) dua kali: $(-x)(-y)=-\bigl(x(-y)\bigr)=-\bigl(-(xy)\bigr)=xy$.
Langkah terakhir memakai $--a=a$, yang juga akibat ketunggalan lawan.

(d) Andaikan $x\neq 0$ dan $xy=0$. Kalikan dengan $1/x$: $(1/x)(xy)=0$, sehingga
$y=0$. Jadi salah satu faktornya nol.
\end{proof}

Tidak satu pun dari keempat pernyataan itu diandaikan; semuanya turun dari hukum
distributif dan keberadaan lawan. Pola ini dipakai sepanjang buku: sebut aturan
sesedikit mungkin, lalu panen akibatnya.

\kamus{lapangan}{field}{Himpunan dengan dua operasi hitung yang aturannya tetap, seperti
kotak alat yang penutupnya selalu rapat: hasil hitungan tidak pernah keluar dari kotak.}

Bab 2 memakai $\mathbb{Q}$ sebagai himpunan terhitung. Di sini ia muncul dalam peran
lain: lapangan yang terurut. Dua sifat itu belum cukup, dan itu isi bagian berikutnya.

\section{Aksioma urutan}
\tingkat{dasar}

Lapangan saja tidak bisa membandingkan dua bilangan. Kita memerlukan relasi ``lebih
kecil dari'' yang menghormati penjumlahan dan perkalian, sebab tanpa itu pertidaksamaan
tidak bisa diolah. Aksioma urutan mengatakan tiga hal: setiap dua bilangan bisa
dibandingkan, perbandingan menular, dan menambah atau mengalikan dengan bilangan positif
tidak merusaknya.

\begin{definisi}[Lapangan terurut]\index{lapangan terurut}
Lapangan $F$ disebut \emph{terurut} jika ada relasi $<$ pada $F$ sehingga untuk setiap
$a,b,c$ berlaku: tepat satu di antara $a<b$, $a=b$, $b<a$ benar; jika $a<b$ dan $b<c$
maka $a<c$; jika $a<b$ maka $a+c<b+c$; dan jika $a<b$ serta $c>0$ maka $ac<bc$.
\end{definisi}

Hampir semua aturan tanda yang kita hafal adalah akibat definisi itu, bukan tambahan.
Bukti berikut menurunkan lima di antaranya sekaligus; polanya sama: tambahkan sesuatu
pada kedua ruas, atau kalikan dengan bilangan yang tandanya diketahui.

\begin{proposisi}
Dalam setiap lapangan terurut berlaku: (a) jika $x>0$ maka $-x<0$, dan sebaliknya;
(b) jika $x>0$ dan $y<z$ maka $xy<xz$; (c) jika $x<0$ dan $y<z$ maka $xy>xz$;
(d) jika $x\neq 0$ maka $x^2>0$, khususnya $1>0$; (e) jika $0<x<y$ maka
$0<1/y<1/x$.
\end{proposisi}

\begin{proof}
(a) Jika $x>0$, tambahkan $-x$ pada kedua ruas: $0>-x$. Sebaliknya, jika $x<0$ maka
$0-x>0$, yaitu $-x>0$.

(b) Karena $z>y$, kita punya $z-y>0$. Kalikan dengan $x>0$: $x(z-y)>0$. Jadi
$xz=x(z-y)+xy>xy$.

(c) Dari (a), $-x>0$. Terapkan (b) pada $-x$ dan $y<z$: $(-x)y<(-x)z$. Menurut
Teorema sebelumnya, ruas kiri adalah $-(xy)$ dan ruas kanan $-(xz)$. Jadi
$-(xy)<-(xz)$, dan dengan (a) lagi $xy>xz$.

(d) Jika $x>0$, aksioma urutan pada $y=x>0$ memberi $x^2>0$. Jika $x<0$ maka $-x>0$,
sehingga $(-x)^2>0$; padahal $(-x)^2=x^2$. Untuk $x=0$ jelas $x^2=0$. Karena
$1=1^2$, kita peroleh $1>0$.

(e) Karena $y>0$ dan $1/y\neq 0$, dari (d) $1/y>0$; sama untuk $1/x$. Kalikan
$x<y$ dengan bilangan positif $(1/x)(1/y)$: hasilnya $1/y<1/x$.
\end{proof}

Dua akibat yang sering dipakai: $a<b$ setara dengan $b-a>0$, dan $a<b$ setara dengan
$-b<-a$; yang kedua menjelaskan mengapa mengalikan dengan bilangan negatif membalik
arah, sebab mengalikan dengan $-1$ mencerminkan garis bilangan terhadap nol. Akibat lain:
jika $0<a<b$ maka $a^2<b^2$, cukup dengan mengalikan $a<b$ dengan $a>0$ lalu dengan
$b>0$ dan menjumlahkannya. Jadi membandingkan kuadrat setara dengan membandingkan
bilangan positif aslinya; setiap kali kita menulis $q^2<2$ untuk $q$ rasional, kita
membandingkan $q$ dengan sesuatu yang belum tentu rasional.

\kamus{lapangan terurut}{ordered field}{Lapangan yang anggotanya bisa diurutkan seperti
anak tangga, dan aturan hitungnya tidak merusak urutan itu.}

Perhatikan apa yang belum kita punya. Kita belum tahu apakah ada bilangan positif yang
lebih kecil dari semua $1/n$, dan belum tahu apakah $\{q\in\mathbb{Q}: q^2<2\}$ punya
atap. Itu bukan soal urutan, melainkan soal kelengkapan.

\section{Nilai mutlak dan ketaksamaan segitiga}
\tingkat{menengah}

Di Bab 2 kita menulis $|x-y|$ sebagai jarak, dan memakainya untuk membandingkan galat,
tanpa membuktikan bahwa lambang itu berperilaku seperti jarak. Sekarang sifatnya
diturunkan dari urutan. Nilai mutlak membuang tanda dan menyisakan jarak ke titik nol;
definisi berikut menulis dua cabang, dan cabang keduanya yang biasa dipakai.

\begin{definisi}[Nilai mutlak]\index{nilai mutlak}
Untuk $x\in\mathbb{R}$, $|x|=x$ jika $x\ge 0$ dan $|x|=-x$ jika $x<0$. Secara setara,
$|x|=\max\{x,-x\}$, dan untuk bilangan real $|x|=\sqrt{x^2}$ dengan akar yang
dibahas di bagian berikutnya.
\end{definisi}

Dari definisi itu langsung: $|x|\ge 0$, $|{-x}|=|x|$, $|x|^2=x^2$, dan $|xy|=|x|\,|y|$.
Yang terakhir terbukti dengan memeriksa empat kombinasi tanda, atau dengan mencatat
bahwa kedua ruas tidak negatif dan kuadratnya sama.

Pertidaksamaan yang memuat nilai mutlak paling mudah diolah setelah diubah menjadi dua
pertidaksamaan biasa. Lemma berikut adalah alasan mengapa perubahan itu sah.

\begin{lemma}
Untuk $c\ge 0$, $|x|\le c$ jika dan hanya jika $-c\le x\le c$.
\end{lemma}

\begin{proof}
Jika $|x|\le c$: untuk $x\ge 0$ kita punya $x=|x|\le c$ dan $-c\le 0\le x$; untuk
$x<0$ kita punya $-x=|x|\le c$, sehingga $x\ge -c$, dan $x<0\le c$. Jadi
$-c\le x\le c$.

Sebaliknya, jika $-c\le x\le c$: untuk $x\ge 0$, $|x|=x\le c$; untuk $x<0$,
$|x|=-x\le c$ karena $x\ge -c$. Dalam kedua kasus $|x|\le c$.
\end{proof}

Ketaksamaan segitiga adalah satu-satunya alat yang kita butuhkan untuk mengendalikan
selisih dua bilangan. Namanya berasal dari gambar: panjang satu sisi segitiga tidak
melebihi jumlah panjang dua sisi lainnya.

\begin{teorema}[Ketaksamaan segitiga]\index{ketaksamaan segitiga}
Untuk semua $x,y\in\mathbb{R}$, $|x+y|\le |x|+|y|$. Selain itu
$\bigl||x|-|y|\bigr|\le |x-y|$.
\end{teorema}

\begin{proof}
Dari definisi, $-|x|\le x\le |x|$ dan $-|y|\le y\le |y|$. Jumlahkan kedua pasangan:
$-(|x|+|y|)\le x+y\le |x|+|y|$. Menurut lemma sebelumnya dengan $c=|x|+|y|$, diperoleh
$|x+y|\le |x|+|y|$.

Untuk bagian kedua, tulis $x=(x-y)+y$ dan terapkan bagian pertama:
$|x|\le |x-y|+|y|$, sehingga $|x|-|y|\le |x-y|$. Dengan menukar peran $x$ dan $y$
diperoleh $|y|-|x|\le |x-y|$. Karena $c\le d$ dan $-c\le d$ setara dengan $|c|\le d$,
maka $\bigl||x|-|y|\bigr|\le |x-y|$.
\end{proof}

Sebagai pemeriksaan, ambil $x=3$ dan $y=-5$: $|x+y|=2$ sedangkan $|x|+|y|=8$, jadi
ketaksamaannya longgar. Untuk $x=3$ dan $y=5$ kedua ruas bernilai $8$ dan ketaksamaan
menjadi persamaan; itu terjadi tepat ketika tandanya sama atau salah satunya nol.

Kegunaan praktisnya: batas pada jarak menjadi batas pada fungsi. Jika $|x-\sqrt2|<\delta$,
maka
\[
|x^2-2|=|x-\sqrt2|\,|x+\sqrt2|<\delta\,(2\sqrt2+\delta),
\]
karena $|x+\sqrt2|\le|x-\sqrt2|+2\sqrt2$. Untuk $\delta=10^{-3}$ batas itu bernilai
$0.002829427$, dan nilai sebenarnya di $x=\sqrt2+10^{-3}$ juga $0.002829427$ sebab
$x^2-2=2\sqrt2\,\delta+\delta^2$ tepat: batas itu ketat. Untuk $\delta=10^{-6}$ batas
yang sama mengecil menjadi $2.828428\times10^{-6}$.

\kamus{nilai mutlak}{absolute value}{Jarak sebuah titik ke nol pada garis bilangan,
tanpa memedulikan sebelah mana titik itu berada.}

Ketaksamaan segitiga hanya memakai aksioma urutan, bukan kelengkapan, jadi ia berlaku
juga di $\mathbb{Q}$ dan tidak bisa membedakan keduanya. Alat pembeda muncul di bagian
berikutnya.

\section{Kelengkapan: batas atas terkecil}
\tingkat{menengah}

Urutan memberi perbandingan, tetapi tidak memberi atap. Himpunan
$A=\{1-1/n: n\in\mathbb{N}\}$ memuat bilangan yang makin dekat ke satu tanpa pernah
sama dengan satu: satu bukan anggotanya, tetapi batasnya paling rapat. Himpunan
$B=\{q\in\mathbb{Q}: q^2<2\}$ terbatas, tetapi tidak punya batas rapat di
$\mathbb{Q}$. Keduanya memperlihatkan bahwa batas longgar dan batas paling rapat perlu
nama berbeda.

\begin{untukanak}
Bayangkan kita harus memasang atap tenda untuk sekelompok anak. Atap terlalu rendah:
ada kepala yang menonjol. Atap terlalu tinggi: tenda boros dan goyah. Yang kita cari
adalah atap paling rendah yang masih membuat semua anak muat. Sebut saja atap pas.

Sekarang ada yang aneh. Di kelompok itu tidak ada anak tertinggi; selalu ada yang
sedikit lebih tinggi. Namun atap pas tetap ada, karena tinggi mereka terbatas: tidak
ada yang setinggi pohon. Atap pas ditentukan oleh seluruh kelompok, bukan oleh satu
anak.

Di bab ini atap pas punya nama: batas atas terkecil. Sifat yang menjamin atap pas
selalu ada disebut kelengkapan. Bilangan pecahan tidak lengkap: ada kelompok yang
tingginya terbatas, tetapi atap pasnya tidak ada di antara bilangan pecahan. Pertanyaan
yang bisa dibawa pulang: kalau tidak ada anak tertinggi, bagaimana caranya membuktikan
bahwa atap pas itu memang yang paling rendah?
\end{untukanak}

Batas atas adalah langit-langit yang memuat semua anggota; batas atas terkecil adalah
langit-langit terendah di antaranya. Definisi berikut menamai keduanya beserta cerminnya
di bawah.

\begin{definisi}[Batas atas dan supremum]\index{batas atas}\index{supremum}
Misalkan $A\subseteq\mathbb{R}$ dan $A\neq\varnothing$. Bilangan $u$ disebut
\emph{batas atas} $A$ jika $a\le u$ untuk setiap $a\in A$. Jika $A$ punya
sekurang-kurangnya satu batas atas, $A$ disebut terbatas di atas. Bilangan $s$ disebut
\emph{supremum} $A$, ditulis $s=\sup A$, jika $s$ batas atas $A$ dan $s\le u$ untuk
setiap batas atas $u$. Secara kembar, $t$ disebut \emph{infimum} $A$, ditulis
$t=\inf A$, jika $t$ batas bawah dan $t\ge v$ untuk setiap batas bawah $v$.
\end{definisi}

Dua syarat pada definisi supremum sering tertukar: yang pertama mengatakan $s$ masih
batas atas, yang kedua bahwa tidak ada batas atas yang lebih kecil. Pada $A=\{1-1/n\}$,
bilangan $1$ memenuhi keduanya sedangkan $2$ hanya yang pertama; himpunan itu juga tidak
punya anggota terbesar, jadi ``maksimum'' dan ``supremum'' bukan sinonim.

\begin{bacakeras}
Dua suara. Suara pertama bertanya, suara kedua menjawab. Bacakan bergantian dengan teman.

``Pagar itu apa?'' ``Pagar adalah batas atas: semua anggota kumpulan ada di bawahnya.''
``Berapa banyak pagarnya?'' ``Banyak. Setiap pagar yang lebih tinggi juga pagar.''
``Lalu yang mana yang kita pakai?'' ``Yang paling rendah. Namanya batas atas terkecil.''
``Kalau pagar terendah itu kebetulan ditempati anggota kumpulan?'' ``Maka ia juga maksimum.
Kalau tidak ditempati, ia tetap batas atas terkecil. Dua nama, dua pekerjaan.''

Ulangi tiga kali, makin cepat: pagar tinggi bukan jawaban, pagar terendah baru jawaban.

Baris yang harus diulang pembaca, sekali, dengan suara jelas: batas atas itu pagar; batas
atas terkecil itu pagar terendah yang masih menahan semua anggota.
\end{bacakeras}

Gambar~\ref{fig:sup-garis} menaruh kedua gagasan itu pada satu garis. Anggota
$A=\{1-1/n\}$ mengerumun ke satu arah tanpa pernah mencapai ujungnya; semua bilangan di
kanan satu adalah batas atas, dan yang paling kiri di antara mereka adalah supremum.

\begin{figure}[htbp]
\centering
\begin{tikzpicture}[x=1.5cm,y=1cm,font=\small]
  \draw[-{Latex[length=2mm]},ink,thick] (-0.35,0) -- (5.5,0);
  \foreach \t in {0,1,2,3,4,5}{
    \draw[ink] (\t,0.07) -- (\t,-0.07);
    \node[below=2pt,ink] at (\t,0) {$\t$};
  }
  \foreach \n in {1,...,6}{
    \fill[accent] ({1-1/\n},1.0) circle (1.7pt);
  }
  \draw[accent,thick,fill=white] (1,1.0) circle (1.7pt);
  \node[anchor=west,accent,font=\footnotesize] at (1.18,1.0) {anggota $A=\{1-1/n\}$};
  \draw[layergrey,dashed] (1,0.95) -- (1,0);
  \node[accent,anchor=west,font=\footnotesize] at (1.12,0.45)
    {$\sup A=1$, dan $1$ bukan anggota $A$};
  \fill[layerteal!14] (1,-0.95) rectangle (5.35,-0.6);
  \draw[layerteal,thick] (1,-0.95) rectangle (5.35,-0.6);
  \node[layerteal!80!black,font=\footnotesize] at (3.2,-0.775) {semua batas atas $A$};
  \draw[accent,line width=1.4pt] (1,-1.02) -- (1,-0.52);
\end{tikzpicture}
\caption{Himpunan $A=\{1-1/n:n\in\mathbb{N}\}$ pada garis bilangan. Setiap $u\ge 1$
adalah batas atas; yang terkecil, $u=1$, adalah $\sup A$. Bilangan $1$ bukan anggota
$A$, jadi $A$ tidak mempunyai maksimum.}
\label{fig:sup-garis}
\end{figure}

\kamus{supremum}{supremum (least upper bound)}{Atap paling rendah yang masih menutupi
seluruh anggota himpunan; atap itu tidak harus bagian dari himpunan.}

Supremum, kalau ada, hanya ada satu; syarat keduanya bisa ditulis dengan epsilon, dan
bentuk itu yang paling sering dipakai dalam bukti.

\begin{lemma}
(i) Jika $A$ mempunyai supremum, maka supremumnya tunggal. (ii) Bilangan $s$ adalah
supremum $A$ jika dan hanya jika $s$ batas atas $A$ dan untuk setiap
$\varepsilon>0$ terdapat $a\in A$ dengan $a>s-\varepsilon$.
\end{lemma}

\begin{proof}
(i) Misalkan $s$ dan $s'$ keduanya supremum. Karena $s$ batas atas dan $s'$ supremum,
$s'\le s$. Karena $s'$ batas atas dan $s$ supremum, $s\le s'$. Jadi $s=s'$.

(ii) Andaikan $s$ batas atas dan syarat epsilon berlaku, tetapi ada batas atas $u<s$.
Ambil $\varepsilon=s-u>0$. Menurut syarat epsilon ada $a\in A$ dengan
$a>s-\varepsilon=u$, bertentangan dengan $u$ batas atas. Sebaliknya, jika $s$ supremum
dan ada $\varepsilon>0$ sehingga $a\le s-\varepsilon$ untuk semua $a\in A$, maka
$s-\varepsilon$ batas atas yang lebih kecil dari $s$, bertentangan dengan $s$ batas
atas terkecil.
\end{proof}

\begin{resep}
Untuk membuktikan bahwa bilangan $s$ adalah supremum himpunan $A$:
\begin{enumerate}
\item Tunjukkan $A\neq\varnothing$ dengan menyebut satu anggotanya.
\item Buktikan $s$ batas atas: ambil $a\in A$ sebarang, tunjukkan $a\le s$.
\item Buktikan tidak ada batas atas yang lebih kecil: ambil $\varepsilon>0$, hasilkan
      $a\in A$ dengan $a>s-\varepsilon$, misalnya $a=1-1/n$ dengan $n>1/\varepsilon$.
\item Simpulkan: $s$ batas atas dan setiap bilangan di bawah $s$ bukan batas atas, jadi
      $s=\sup A$.
\end{enumerate}
Peringatan urutan: langkah 2 harus selesai sebelum langkah 3. Memilih $a>s-\varepsilon$
lebih dulu tidak membuktikan apa pun, sebab $s$ belum diketahui batas atas. Pada
$A=\{1-1/n\}$, bilangan $1/2$ lulus langkah 3 tetapi gagal langkah 2.
\end{resep}

Aksioma berikut membedakan $\mathbb{R}$ dari $\mathbb{Q}$. Ia tidak bisa dibuktikan dari
hukum lapangan dan urutan, sebab ia pernyataan tentang kelengkapan garis.

\begin{teorema}[Sifat batas atas terkecil]\index{aksioma kelengkapan}
Setiap himpunan tak kosong $A\subseteq\mathbb{R}$ yang terbatas di atas mempunyai
supremum di $\mathbb{R}$.
\end{teorema}

\begin{proof}[Bukti]
Pernyataan ini diterima sebagai aksioma kelengkapan. Rumusan lain yang setara: irisan
selang bersarang, dan kekonvergenan barisan Cauchy. Kita memilih supremum karena paling
langsung dipakai; pembangunan $\mathbb{R}$ dari $\mathbb{Q}$ yang membuktikannya ada
pada Apendiks Bab 1 Rudin.
\end{proof}

Infimum tidak perlu diaksiomakan: ia supremum yang dilihat di cermin.

\begin{proposisi}
Jika $A$ tak kosong dan terbatas di bawah, maka $\inf A=-\sup(-A)$, dengan
$-A=\{-a: a\in A\}$.
\end{proposisi}

\begin{proof}
Himpunan $-A$ tak kosong dan terbatas di atas, jadi $s=\sup(-A)$ ada. Untuk setiap
$a\in A$ berlaku $-a\le s$, sehingga $a\ge -s$: jadi $-s$ batas bawah $A$. Sekarang
ambil $t>-s$. Maka $-t<s$, sehingga $-t$ bukan batas atas $-A$; ada $a\in A$ dengan
$-a>-t$, yaitu $a<t$. Jadi $t$ bukan batas bawah $A$. Karena itu $-s$ batas bawah
terbesar.
\end{proof}

Akibat pertama kelengkapan: bilangan asli tidak punya langit-langit. Kedengarannya
sepele, tetapi tanpa kelengkapan ia tidak bisa dibuktikan.

\begin{teorema}[Sifat Archimedes]\index{sifat archimedes}
Untuk setiap $x>0$ dan setiap $y\in\mathbb{R}$ terdapat $n\in\mathbb{N}$ dengan
$nx>y$.
\end{teorema}

\begin{proof}
Andaikan sebaliknya: ada $x>0$ dan $y$ sehingga $nx\le y$ untuk semua $n$. Maka himpunan
$S=\{nx: n\in\mathbb{N}\}$ tak kosong dan terbatas di atas oleh $y$, sehingga
$\alpha=\sup S$ ada. Karena $x>0$, bilangan $\alpha-x$ lebih kecil dari $\alpha$ dan
karena itu bukan batas atas $S$. Jadi ada $m$ dengan $mx>\alpha-x$, yang memberi
$(m+1)x>\alpha$. Padahal $(m+1)x\in S$, dan $\alpha$ batas atas $S$. Kontradiksi.
\end{proof}

Versi resiprokalnya yang paling sering dipakai: untuk setiap $\varepsilon>0$ ada
$n$ dengan $1/n<\varepsilon$, cukup dengan menerapkan teorema pada $x=\varepsilon$ dan
$y=1$. Kalimat ``untuk $n$ cukup besar'' adalah singkatan dari akibat itu, dan sekarang
kita tahu ia sah.

Sifat itu juga memusnahkan bilangan positif yang sangat kecil: tidak ada $\delta>0$
yang lebih kecil dari $1/n$ untuk semua $n$, sebab $1/\delta$ akan menjadi batas atas
$\mathbb{N}$. Karena itu analisis real tidak memerlukan ``infinitesimal''.

Akibat kedua: setiap selang terbuka, sekecil apa pun, memuat bilangan rasional.

\begin{teorema}[Kerapatan $\mathbb{Q}$]\index{kerapatan}
Jika $x<y$ bilangan real, maka ada $p\in\mathbb{Q}$ dengan $x<p<y$.
\end{teorema}

\begin{proof}
Karena $y-x>0$, Sifat Archimedes memberi $n\in\mathbb{N}$ dengan $n(y-x)>1$. Terapkan
Archimedes lagi untuk memperoleh bilangan bulat $m_1>nx$ dan $m_2>-nx$; jadi
$-m_2<nx<m_1$. Karena $nx$ terkurung di antara dua bilangan bulat, ada bilangan bulat
$m$ dengan $m-1\le nx<m$. Maka $nx<m$, dan dari $m-1\le nx$ diperoleh
$m\le nx+1<ny$. Karena $n>0$, pembagian memberi $x<m/n<y$.
\end{proof}

Bukti itu konstruktif: ia memberi resep untuk menemukan $p$, yaitu memilih $n$ cukup
besar lalu membulatkan $nx$ ke atas. Ia tidak memberi penyebut yang kecil, jadi
aproksimasi efisien masih menuntut gagasan lain.

\section{Mengapa bilangan rasional belum cukup}
\tingkat{lanjut}

Program di pembuka gagal karena himpunan yang ia apit tidak punya atap di
$\mathbb{Q}$. Mari kita buktikan, bukan hanya menunjukkan angkanya. Ambil
$A=\{q\in\mathbb{Q}: q^2<2\}$. Himpunan ini tidak kosong, sebab $1\in A$, dan terbatas
di atas di dalam $\mathbb{Q}$: jika $q\ge 3/2$ maka $q^2\ge 9/4>2$, jadi $3/2$ batas
atas.

Selang-selang yang dipertahankan program itu digambar sebagai tumpukan batang yang makin
pendek (Gambar~\ref{fig:selang-bersarang}). Setiap batang masih memuat akar dua dan
panjangnya menjadi separuh, sementara semua ujungnya tetap rasional; batang-batang itu
menyempit ke satu titik tanpa pernah menyentuhnya.

\begin{figure}[htbp]
\centering
\begin{tikzpicture}[x=9cm,y=1cm,font=\small]
  \foreach \y/\a/\b/\lab/\pan in {%
    0/1/2/{$[1,\;2]$}/{$1$},
    0.55/1/1.5/{$[1,\;\tfrac32]$}/{$1/2$},
    1.1/1.25/1.5/{$[\tfrac54,\;\tfrac32]$}/{$1/4$},
    1.65/1.375/1.5/{$[\tfrac{11}{8},\;\tfrac32]$}/{$1/8$},
    2.2/1.375/1.4375/{$[\tfrac{11}{8},\;\tfrac{23}{16}]$}/{$1/16$},
    2.75/1.40625/1.4375/{$[\tfrac{45}{32},\;\tfrac{23}{16}]$}/{$1/32$}}{
    \fill[layerindigo!12] (\a,\y) rectangle (\b,\y+0.3);
    \draw[layerindigo,thick] (\a,\y) rectangle (\b,\y+0.3);
    \node[layerindigo,anchor=east,font=\footnotesize] at (\a-0.018,\y+0.15) {\lab};
    \node[layergrey,anchor=west,font=\footnotesize] at (\b+0.015,\y+0.15) {\pan};
  }
  \draw[accent,dashed,thick] (1.41421,-0.25) -- (1.41421,3.15);
  \node[accent,font=\footnotesize,above] at (1.41421,3.15) {$\sqrt2$};
  \draw[-{Latex[length=1.6mm]},layergrey] (1.0,-0.12) -- (2.0,-0.12);
  \node[layergrey,font=\footnotesize,below] at (1.5,-0.14) {panjang $=1$};
\end{tikzpicture}
\caption{Selang yang dipertahankan program pembuka: $[1,2]$, $[1,\tfrac32]$,
$[\tfrac54,\tfrac32]$, $[\tfrac{11}{8},\tfrac32]$, $[\tfrac{11}{8},\tfrac{23}{16}]$,
$[\tfrac{45}{32},\tfrac{23}{16}]$. Panjangnya turun dari $1$ ke $1/32$; diteruskan tanpa
henti, irisannya menyusut ke satu titik, yaitu $\sqrt2$, yang bukan bilangan rasional.}
\label{fig:selang-bersarang}
\end{figure}

Langkah pertama: tidak ada anggota terbesar. Setiap rasional dengan $q^2<2$ masih bisa
digeser sedikit ke atas tanpa melewati dua, dan besar geserannya punya rumus.

\begin{lemma}
Jika $q\in\mathbb{Q}$, $q>0$, dan $q^2<2$, maka terdapat $h\in\mathbb{Q}$, $h>0$,
dengan $(q+h)^2<2$. Khususnya $q+h\in A$ dan $q+h>q$.
\end{lemma}

\begin{proof}
Ambil $h=\dfrac{2-q^2}{2q+2}$. Karena $q^2<2$ dan $2q+2>0$, maka $h>0$ dan $h$
rasional. Juga $h<1$, sebab $2-q^2<2q+2$ setara dengan $q^2+2q>0$, yang benar untuk
$q>0$. Maka $h^2\le h$, sehingga
\[
(q+h)^2=q^2+2qh+h^2\le q^2+h(2q+1)
=q^2+(2-q^2)\frac{2q+1}{2q+2}<q^2+(2-q^2)=2,
\]
karena $\frac{2q+1}{2q+2}<1$.
\end{proof}

Rumus itu bisa diperiksa dengan angka. Untuk $q=1393/985$ berlaku $2-q^2=1/970225$,
sehingga $h=1/4684660$ dan
\[
q+h=\frac{6625109}{4684660}=1.414213411432206\ldots,
\]
yang kuadratnya $1.999999573074719\ldots<2$. Jadi $1393/985$ bukan anggota terbesar,
meskipun galat kuadratnya sudah sekitar satu per sejuta. Prosedur ini bisa diulang tanpa
henti: tidak ada langkah terakhir.

\begin{tangan}
Bahan: tali 1,5 m, enam lakban, meteran.

Tempel lakban di lantai tegak lurus dinding, pada jarak $60/n$ cm untuk $n=1,\ldots,6$:
60, 30, 20, 15, 12, dan 10 cm. Kamu baru menempel enam, tetapi aturannya tidak berhenti.
Bentangkan tali melintang sebagai pagar, lalu geser ke arah dinding. Setiap kali
berhenti, periksa satu hal: apakah semua lakban ada di sisi tali yang sama dengan
dinding? Selama ya, tali masih pagar.

Tali bisa didekatkan ke dinding tanpa batas dan tetap pagar: 5 cm, 1 cm, 1 mm. Yang
menghentikannya bukan lakban, melainkan dinding. Kumpulan lakbannya tidak punya anggota
terkecil, dan tidak satu pun menempel di dinding. Pagar terendah tetap ada---dinding---
dan ia bukan anggota kumpulan. Itulah supremum.
\end{tangan}

Bukan hanya maksimumnya yang tidak ada. Batas atas terkecilnya pun tidak ada di dalam
bilangan rasional, dan buktinya memerlukan satu fakta tentang akar dua.

\begin{proposisi}
Himpunan $A=\{q\in\mathbb{Q}: q^2<2\}$ tidak mempunyai supremum di $\mathbb{Q}$.
\end{proposisi}

\begin{proof}
Andaikan $p\in\mathbb{Q}$ supremum $A$. Karena $1\in A$, maka $p\ge 1>0$. Kita
tunjukkan $p^2=2$, yang mustahil karena tidak ada bilangan rasional yang kuadratnya dua.

Jika $p^2<2$, lemma sebelumnya memberi $q=p+h\in A$ dengan $q>p$, bertentangan dengan
$p$ batas atas. Jika $p^2>2$, ambil $k=(p^2-2)/(2p)>0$, yang rasional dan lebih kecil
dari $p$ sebab $p^2-2<2p^2$. Untuk setiap $q\in A$,
\[
(p-k)^2=p^2-2pk+k^2>p^2-2pk=2>q^2,
\]
jadi $q<p-k$ untuk $q>0$, dan untuk $q\le 0$ hal itu jelas karena $p-k>0$: bilangan
$p-k$ masih batas atas $A$, tetapi lebih kecil dari $p$.
Bertentangan dengan $p$ batas atas terkecil. Jadi $p^2=2$.

Sekarang $p=m/n$ dengan $m,n$ bulat tanpa faktor bersama. Dari $m^2=2n^2$, $m^2$ genap,
sehingga $m$ genap; tulis $m=2r$. Maka $4r^2=2n^2$, yaitu $n^2=2r^2$, sehingga $n$
juga genap. Keduanya genap, bertentangan dengan pilihan $m,n$.
\end{proof}

Di dalam $\mathbb{R}$ hal itu tidak terjadi. Himpunan yang sama, kali ini dipandang
sebagai himpunan bilangan real, terbatas di atas oleh $2$, jadi menurut aksioma
kelengkapan ia punya supremum. Supremum itu adalah bilangan yang kita namai $\sqrt{2}$.
Bukti berikut menutup dua kemungkinan yang tersisa, dan bentuknya meniru bukti lemma
tadi dengan $q$ diganti bilangan real.

\begin{teorema}[Keberadaan akar dua]\index{akar dua}
Terdapat tepat satu $s>0$ dengan $s^2=2$.
\end{teorema}

\begin{proof}
Misalkan $E=\{t\in\mathbb{R}: t>0,\ t^2<2\}$. Himpunan ini tak kosong karena $1\in E$,
dan terbatas di atas oleh $2$: jika $t>2$ maka $t^2>4>2$, jadi $t\notin E$. Menurut
sifat batas atas terkecil, $s=\sup E$ ada, dan $s\ge 1>0$ karena $1\in E$.

Andaikan $s^2<2$. Ambil $h=\frac12\min\{1,(2-s^2)/(2s+1)\}$. Maka $0<h<1$ dan
$h(2s+1)<2-s^2$: jika nilai minimum yang pertama, $h=\frac12$ dan
$\frac12(2s+1)\le\frac12(2-s^2)<2-s^2$; jika yang kedua, $h(2s+1)=\frac12(2-s^2)$.
Karena $h<1$, kita punya $h^2<h$, sehingga
\[
(s+h)^2=s^2+2sh+h^2<s^2+h(2s+1)<s^2+(2-s^2)=2 .
\]
Jadi $s+h\in E$ padahal $s+h>s$, bertentangan dengan $s$ batas atas $E$. Karena itu
$s^2\ge 2$.

Andaikan $s^2>2$. Ambil $k=(s^2-2)/(2s)>0$; maka $k<s$ karena $s^2-2<2s^2$. Untuk
setiap $t\in E$ dengan $t\ge s-k$ berlaku
\[
t^2\ge(s-k)^2=s^2-2sk+k^2>s^2-2sk=2,
\]
sehingga $t\notin E$. Jadi $s-k$ batas atas $E$ yang lebih kecil dari $s$, bertentangan
dengan $s$ batas atas terkecil. Karena $s^2<2$ dan $s^2>2$ keduanya mustahil, $s^2=2$.

Ketunggalan: jika $s,t>0$ dan $s^2=t^2=2$, maka $(s-t)(s+t)=0$. Karena $s+t>0$,
$s=t$.
\end{proof}

Bilangan $s$ irasional, menurut argumen paritas tadi: lubang di $\mathbb{Q}$ nyata, dan
$\mathbb{R}$ menutupnya dengan tepat satu bilangan. Yang menutupnya bukan urutan dan
bukan hukum lapangan, melainkan kelengkapan. Program pembuka tidak gagal karena kurang
teliti, melainkan karena sistem bilangannya tidak memuat objek yang dicarinya.

Satu hal perlu dinyatakan jujur: kita belum membangun $\mathbb{R}$ dari $\mathbb{Q}$,
hanya mengandaikan $\mathbb{R}$ memenuhi tiga kelompok aksioma. Pembangunannya dikenal.
Dedekind memotong $\mathbb{Q}$ menjadi dua: sebuah \emph{potongan} adalah himpunan
bagian $\alpha\subset\mathbb{Q}$ yang tidak kosong, tidak seluruh $\mathbb{Q}$, turun
tertutup, dan tanpa anggota terbesar. Sifat terakhir persis lemma kita tadi. Cantor
memakai kelas ekuivalen barisan Cauchy rasional; keduanya terbit pada 1872. Rinciannya
ada pada Apendiks Bab 1 Rudin.

Pembangunan itu menjelaskan kedudukan desimal. Untuk $x>0$, ambil $n_0$ bilangan bulat
terbesar dengan $n_0\le x$, lalu $n_1$ angka terbesar dengan $n_0+n_1/10\le x$, lalu
$n_2$ dengan $n_0+n_1/10+n_2/100\le x$, dan seterusnya. Pemotongan itu terbatas di atas
oleh $x$, jadi supremumnya ada dan sama dengan $x$. Desimal bukan bahan dasar bilangan
real, melainkan cara menamai supremum; kelengkapan menjamin namanya ada, bukan bahwa
namanya bisa ditulis dengan pecahan.

\kamus{aksioma kelengkapan}{completeness axiom}{Janji bahwa garis bilangan tidak
berlubang: setiap kumpulan yang terbatas ke atas selalu punya atap paling rendah.}

Keberadaan $\sqrt2$ membuat soal komputasi masuk akal: $s$ bisa dihitung tanpa tahu
digitnya lebih dulu. Iterasi Newton, $x_{k+1}=\frac12(x_k+2/x_k)$, tetap rasional bila
$x_k$ rasional, sehingga seluruh proses berjalan di dalam $\mathbb{Q}$. Dari $s^2=2$
diperoleh
\[
x_{k+1}-s=\frac{(x_k-s)^2}{2x_k},\qquad e_{k+1}=e_k^2\cdot\frac{s}{2x_k},
\]
dengan $e_k=|x_k-s|/s$. Ruas kanan bentuk pertama tidak negatif, jadi setiap langkah
tetap di atas $s$; di sana faktor $\frac{s}{2x_k}$ tidak melebihi $1/2$, sehingga
$e_{k+1}\le e_k^2/2$ dan jumlah digit benar kira-kira berlipat dua setiap langkah.

Satu peringatan: kita belum membuktikan $x_k\to s$. Untuk itu diperlukan alat Bab 4.
Yang bisa dikatakan sekarang: karena $e_1<0.061$, dari $e_{k+1}\le e_k^2/2$ dapat
diturunkan $e_k\le 2^{-(2^k-1)}$ untuk $k\ge 1$, dan itu cukup untuk merancang uji
numerik yang jujur.

\section{Interval, titik akumulasi, dan himpunan terbuka}
\tingkat{lanjut}

Dengan supremum di tangan, kita bisa memberi nama pada potongan garis bilangan, seperti
aturan finis lomba lari: garisnya sama, yang berbeda hanya apakah garis itu ikut
dihitung. Nama-nama itu dipakai lagi di Bab 11 dan Bab 13, jadi kita ambil yang perlu
saja tanpa membangun topologi secara penuh.

Interval adalah potongan garis yang tidak berlubang. Definisi berikut menuliskan
sembilan bentuk yang mungkin; semuanya ditentukan oleh dua ujung dan oleh keputusan
apakah ujung itu ikut dihitung.

\begin{definisi}[Interval]
Misalkan $a<b$ bilangan real. Himpunan $(a,b)=\{x: a<x<b\}$ disebut interval terbuka,
$[a,b]=\{x: a\le x\le b\}$ interval tertutup, dan $[a,b)$, $(a,b]$ adalah bentuk
campuran. Himpunan $(a,\infty)=\{x: x>a\}$, $[a,\infty)$, $(-\infty,b)$,
$(-\infty,b]$, dan $\mathbb{R}$ sendiri disebut interval tak terbatas. Kesembilan
bentuk itu disebut interval.
\end{definisi}

Ciri yang menandai semuanya: jika dua titik ada di dalam interval, seluruh ruas di
antaranya juga di dalam, sehingga interval disebut tidak berlubang. Daftar tadi tidak
dibuktikan lengkap; klasifikasinya hanya memeriksa batas atas, batas bawah, dan apakah
batas itu termasuk. Titik ``dekat'' dinyatakan dengan persekitaran, yaitu interval
terbuka kecil di sekitar satu titik.

\begin{definisi}[Persekitaran, titik dalam, terbuka, tertutup]\index{himpunan terbuka}\index{himpunan tertutup}\index{titik akumulasi}
Untuk $\varepsilon>0$, persekitaran titik $x$ adalah
$V_\varepsilon(x)=(x-\varepsilon,x+\varepsilon)$. Titik $x$ disebut titik dalam
$S\subseteq\mathbb{R}$ jika $V_\varepsilon(x)\subseteq S$ untuk suatu $\varepsilon>0$;
$S$ disebut terbuka jika setiap anggotanya titik dalam. Titik $x$ disebut titik
akumulasi $S$ jika setiap $V_\varepsilon(x)$ memuat anggota $S$ yang berbeda dari $x$.
Himpunan $S$ disebut tertutup jika komplemennya terbuka.
\end{definisi}

Beberapa contoh menata intuisi. Interval $(0,1)$ terbuka dan tidak tertutup, sebab
setiap persekitaran $0$ memotongnya. Interval $[0,1]$ tertutup dan tidak terbuka,
sedangkan $[0,1)$ bukan keduanya. Himpunan $\mathbb{Q}$ juga bukan keduanya, sebab
setiap persekitaran memuat rasional dan irasional. Himpunan $\{1/n:n\in\mathbb{N}\}$
tidak tertutup: $0$ ada di komplemennya, tetapi setiap persekitaran $0$ memuat $1/n$.
Himpunan $\mathbb{Z}$ tertutup, sebab komplemennya gabungan interval terbuka $(k,k+1)$.

Jadi ``tertutup'' bukan negasi dari ``terbuka'': $\mathbb{R}$ dan $\varnothing$
keduanya terbuka sekaligus tertutup, dan komplemen himpunan terbuka selalu tertutup.
Kerapatan $\mathbb{Q}$ membuat contoh itu lebih tajam. Setiap persekitaran memuat
rasional dan irasional, jadi tidak ada titik yang menjadi titik dalam $\mathbb{Q}$
maupun komplemennya: himpunan itu padat sekaligus berlubang di mana-mana. Karena itu
kelengkapan tidak bisa diganti dengan kerapatan.

Persekitaran sendiri terbuka, dan itu membuat definisi titik dalam bisa dipakai
berulang.

\begin{proposisi}
(i) Setiap persekitaran $V_\varepsilon(x)$ terbuka. (ii) Gabungan sebarang keluarga
himpunan terbuka adalah terbuka. (iii) Irisan berhingga himpunan terbuka adalah
terbuka.
\end{proposisi}

\begin{proof}
(i) Ambil $y\in V_\varepsilon(x)$, jadi $|y-x|<\varepsilon$. Tetapkan
$\delta=\varepsilon-|y-x|>0$. Untuk setiap $z\in V_\delta(y)$,
$|z-x|\le|z-y|+|y-x|<\delta+|y-x|=\varepsilon$ menurut ketaksamaan segitiga. Jadi
$V_\delta(y)\subseteq V_\varepsilon(x)$.

(ii) Jika $x$ berada di gabungan, ia berada di salah satu anggotanya, dan karena
anggota itu terbuka ada persekitaran $x$ yang termuat di dalamnya, sehingga juga
termuat di dalam gabungan.

(iii) Jika $x$ berada di irisan berhingga $S_1\cap\cdots\cap S_n$, untuk setiap $i$ ada
$\varepsilon_i>0$ dengan $V_{\varepsilon_i}(x)\subseteq S_i$. Ambil
$\varepsilon=\min\{\varepsilon_1,\ldots,\varepsilon_n\}>0$. Maka
$V_\varepsilon(x)$ termuat di setiap $S_i$, jadi termuat di irisannya.
\end{proof}

Butir (iii) hanya untuk irisan berhingga: $\bigcap_n(-1/n,1/n)=\{0\}$ tidak terbuka,
sebab untuk $x\neq0$ Sifat Archimedes memberi $n$ dengan $1/n<|x|$.

\section{Contoh hitung}
\tingkat{menengah}

Tiga perhitungan yang dijanjikan dijalankan di sini: iterasi Newton, aproksimasi
rasional sampai galat di bawah $10^{-6}$, dan pemeriksaan bahwa $\{q:q^2<2\}$ tidak
punya maksimum. Semua angkanya dihitung dengan pecahan eksak.

Mulai dari $x_0=1$, iterasi $x_{k+1}=(x_k+2/x_k)/2$ memberi pecahan berikut; penyebutnya
tumbuh, dan itu harga yang dibayar untuk ketepatan.

\begin{center}\small
\begin{tabular}{rlll}
\hline
$k$ & $x_k$ & desimal & galat relatif \\
\hline
0 & $1$ & $1$ & $2.92893\times10^{-1}$ \\
1 & $3/2$ & $1.5$ & $6.06602\times10^{-2}$ \\
2 & $17/12$ & $1.41666666667$ & $1.73461\times10^{-3}$ \\
3 & $577/408$ & $1.41421568627$ & $1.50183\times10^{-6}$ \\
4 & $665857/470832$ & $1.41421356237$ & $1.12774\times10^{-12}$ \\
5 & $886731088897/627013566048$ & $1.41421356237$ & $6.35896\times10^{-25}$ \\
\hline
\end{tabular}
\end{center}

Galat relatif dihitung terhadap $\sqrt{2}=1.414213562373095\ldots$. Faktor
penurunannya naik sendiri: dari langkah 0 ke 1 sekitar $4.8$, lalu $35$, $1155$, dan
$1.3\times10^{6}$ --- itulah tanda pengkuadratan. Nisbah $e_{k+1}/e_k^2$ bernilai
$0.7071$, $0.4714$, $0.4991$, $0.49999925$, dan $0.5$, sesuai rumus
$e_{k+1}=e_k^2\cdot\frac{s}{2x_k}$ karena $x_k$ mendekat ke $s$ dari atas. Pada langkah
4 dan 5 desimalnya sama sampai sebelas angka; angka berikutnya tidak muncul karena kita
mencetak dua belas angka penting, bukan karena galatnya berhenti turun. Pecahan langkah
5 punya dua belas digit dengan galat $6.36\times10^{-25}$, dan satu angka tambahan hanya
menuntut kira-kira satu langkah lagi.

Cara kedua memakai aproksimasi rasional yang lebih hemat penyebut, yaitu pecahan lanjar
$\sqrt2$:
\[
1,\quad \frac32,\quad \frac75,\quad \frac{17}{12},\quad \frac{41}{29},\quad \ldots
\]
dengan $p_{k+1}=p_k+2q_k$ dan $q_{k+1}=p_k+q_k$ mulai dari $(p_0,q_0)=(1,1)$. Tabel
berikut memuat galat mutlaknya.

\begin{center}\small
\begin{tabular}{rlll}
\hline
$p/q$ & desimal & galat mutlak & $p^2-2q^2$ \\
\hline
$1/1$ & $1$ & $4.14214\times10^{-1}$ & $-1$ \\
$3/2$ & $1.5$ & $8.57864\times10^{-2}$ & $1$ \\
$7/5$ & $1.4$ & $1.42136\times10^{-2}$ & $-1$ \\
$17/12$ & $1.41666666667$ & $2.45310\times10^{-3}$ & $1$ \\
$41/29$ & $1.41379310345$ & $4.20459\times10^{-4}$ & $-1$ \\
$99/70$ & $1.41428571429$ & $7.21519\times10^{-5}$ & $1$ \\
$239/169$ & $1.41420118343$ & $1.23789\times10^{-5}$ & $-1$ \\
$577/408$ & $1.41421568627$ & $2.12390\times10^{-6}$ & $1$ \\
$1393/985$ & $1.41421319797$ & $3.64404\times10^{-7}$ & $-1$ \\
$3363/2378$ & $1.41421362489$ & $6.25218\times10^{-8}$ & $1$ \\
\hline
\end{tabular}
\end{center}

Aproksimasi pertama yang galatnya di bawah $10^{-6}$ adalah $1393/985$ dengan galat
$3.64\times10^{-7}$; angka sebelumnya, $577/408$, masih $2.12\times10^{-6}$. Kolom
terakhir menjelaskan sebabnya: karena $p^2-2q^2=\pm1$,
\[
\left|\frac pq-\sqrt2\right|
=\frac{|p^2-2q^2|}{q^2\left(\frac pq+\sqrt2\right)}
<\frac{1}{2\sqrt2\,q^2},
\]
sehingga galat turun seperti kebalikan kuadrat penyebut. Untuk $q=985$ batas itu
$3.644\times10^{-7}$ dan galat sebenarnya $3.64404\times10^{-7}$: hampir tepat, karena
$\frac pq$ sudah sangat dekat ke $\sqrt2$. Penyebut 985 jauh lebih kecil daripada
$627013566048$ pada langkah 5 Newton, meskipun galatnya hanya sekitar sejuta kali lebih
besar --- itu beda metode yang mengejar galat kuadrat dan yang mengejar galat rasional
terkecil.

Gambar~\ref{fig:tangga-akar} menampilkan kolom galat itu sebagai tangga menurun. Tanda
$p^2-2q^2$ bergantian, sehingga pecahannya melompat dari bawah ke atas akar dua, dan
nisbah dua galat berurutan menuju $5{,}8284=(1+\sqrt2)^2$.

\begin{figure}[htbp]
\centering
\begin{tikzpicture}[x=1.02cm,y=1cm,font=\small]
  \foreach \lv in {1,...,7}{
    \draw[layergrey!20] (0.55,\lv) -- (10.75,\lv);
    \node[layergrey,anchor=east,font=\scriptsize] at (0.45,\lv) {$10^{-\lv}$};
  }
  \node[layergrey,anchor=east,font=\scriptsize] at (0.45,0) {$10^{0}$};
  \draw[-{Latex[length=2mm]},ink,thick] (0.55,0) -- (10.95,0);
  \draw[-{Latex[length=2mm]},ink,thick] (0.55,0) -- (0.55,7.85);
  \foreach \k/\lab in {%
    1/{$1/1$},2/{$3/2$},3/{$7/5$},4/{$17/12$},5/{$41/29$},
    6/{$99/70$},7/{$239/169$},8/{$577/408$},9/{$1393/985$},10/{$3363/2378$}}{
    \draw[ink] (\k,0.07) -- (\k,-0.07);
    \node[rotate=32,anchor=north east,font=\scriptsize,layergrey] at (\k,-0.12) {\lab};
  }
  \foreach \k/\v in {1/0.3828,3/1.8473,5/3.3763,7/4.9073,9/6.4384}{
    \fill[accent] (\k,\v) circle (1.8pt);
  }
  \foreach \k/\v in {2/1.0666,4/2.6103,6/4.1418,8/5.6729,10/7.2040}{
    \fill[layerteal] (\k,\v) circle (1.8pt);
  }
  \draw[layergrey,dashed,thick] (0.55,6) -- (10.75,6);
  \node[layergrey,font=\scriptsize,anchor=south east] at (10.7,6.02) {ambang $10^{-6}$};
  \fill[accent] (0.95,7.55) circle (1.8pt);
  \node[right,font=\scriptsize] at (1.05,7.55) {di bawah $\sqrt2$};
  \fill[layerteal] (4.35,7.55) circle (1.8pt);
  \node[right,font=\scriptsize] at (4.45,7.55) {di atas $\sqrt2$};
  \node[font=\scriptsize,anchor=north] at (5.7,-1.18) {pecahan berurutan};
  \node[font=\scriptsize,rotate=90,anchor=south] at (-0.35,3.9) {$|p/q-\sqrt2|$};
\end{tikzpicture}
\caption{Tangga aproksimasi rasional: galat $|p/q-\sqrt2|$ bagi sepuluh pecahan
berurutan, skala logaritmik. Ambang $10^{-6}$ pertama dilewati $1393/985$, dengan galat
$3{,}644\times10^{-7}$; sesudahnya tiap anak tangga turun dengan faktor $5{,}8284$.}
\label{fig:tangga-akar}
\end{figure}

Pemeriksaan ketiga menyangkut ketiadaan maksimum. Satu langkah lagi dari $q$ sebelumnya
memberi
\[
q'=\frac{149856914148961}{105964844887080}=1.41421349985133266\ldots,
\]
dan kuadratnya masih di bawah dua. Kedua bilangan itu rasional, keduanya di dalam $A$,
dan masing-masing lebih besar dari yang sebelumnya: tidak ada alasan proses ini berhenti.

\begin{tantangan}
Tinjau $S=\{q\in\mathbb{Q}: q\ge 0,\ q^2\le 2\}$. Sekilas himpunan ini punya anggota
terbesar: bilangan yang kuadratnya tepat dua, sebab tanda $\le$ seolah memasukkan batas
itu ke dalam himpunan. Buktikan bahwa $S$ tidak punya maksimum, lalu tentukan $\sup S$
di $\mathbb{R}$ dan jelaskan mengapa jawabannya bukan anggota $S$.

Petunjuk 1: anggota $S$ harus rasional. Periksa dulu apakah calon terbesarmu boleh
menjadi anggota himpunan.

Petunjuk 2: untuk setiap $q\in S$ dengan $q>0$ berlaku $q^2<2$ secara tegas, sebab
$q^2=2$ mustahil bagi rasional. Jadi lemma geseran $h=(2-q^2)/(2q+2)$ bekerja dan
menghasilkan anggota $S$ yang lebih besar; batas yang tak pernah tercapai itu $\sqrt2$.

Kalau mentok, kembali ke bagian ``Mengapa bilangan rasional belum cukup'' dan ke lemma
tentang geseran.
\end{tantangan}

Terakhir, Sifat Archimedes pada angka. Untuk $\varepsilon=10^{-6}$ kita perlu
$n>10^{6}$. Pada $n=10^{6}$ resiprokalnya tepat $10^{-6}$, yang tidak lebih kecil dari
$\varepsilon$; pada $n=10^{6}+1$ resiprokalnya $9.99999\times10^{-7}$, baru memenuhi.
Perbedaan satu pada $n$ menentukan sah tidaknya pilihan, dan itulah sebabnya
pertidaksamaan tegas harus dihormati.

\kertas{Untuk $x=1393/985$, hitung $h=(2-x^2)/(2x+2)$ dan periksa dengan tangan bahwa
$(x+h)^2<2$. Lalu hitung galat relatif $|x-\sqrt2|/\sqrt2$ dan bandingkan dengan
$1/(2\sqrt2\,x^2)$. Jelaskan mengapa keduanya hampir sama, dan mengapa fakta itu tidak
bertentangan dengan ketiadaan maksimum pada $A$.}

\section{Lab komputasi}
\tingkat{lanjut}

Kode berikut dijalankan dengan Python 3 dan \texttt{mpmath} untuk aritmetika presisi
tinggi; iterasi Newton memakai \texttt{Fraction}, jadi tidak ada pembulatan sebelum
pencetakan.

\begin{lab}
\begin{verbatim}
from fractions import Fraction as F
import mpmath as mp
mp.mp.dps = 40
akar2 = mp.sqrt(2)

print("Newton untuk akar dua, x0 = 1")
x = F(1)
for k in range(6):
    v = mp.mpf(x.numerator) / x.denominator
    print(k, f"{x.numerator}/{x.denominator}", mp.nstr(v, 12),
          mp.nstr(abs(v - akar2) / akar2, 6))
    x = (x + F(2) / x) / 2

print()
print("Aproksimasi rasional berurutan, galat < 1e-6 ?")
p, q = 1, 1
for k in range(10):
    v = mp.mpf(p) / q
    print(f"{p}/{q}", mp.nstr(v, 12), mp.nstr(abs(v - akar2), 6),
          abs(v - akar2) < mp.mpf(10) ** -6)
    p, q = p + 2 * q, p + q

print()
print("Uji: {q rasional : q^2 < 2} tidak punya maksimum")
q = F(1393, 985)
for k in range(2):
    h = (2 - q * q) / (2 * q + 2)
    q = q + h
    print(q, mp.nstr(mp.mpf(q.numerator) / q.denominator, 18), q * q < 2)
\end{verbatim}
\end{lab}

Keluaran yang diperoleh, apa adanya:

\begin{verbatim}
Newton untuk akar dua, x0 = 1
0 1/1 1.0 0.292893
1 3/2 1.5 0.0606602
2 17/12 1.41666666667 0.00173461
3 577/408 1.41421568627 1.50183e-6
4 665857/470832 1.41421356237 1.12774e-12
5 886731088897/627013566048 1.41421356237 6.35896e-25

Aproksimasi rasional berurutan, galat < 1e-6 ?
1/1 1.0 0.414214 False
3/2 1.5 0.0857864 False
7/5 1.4 0.0142136 False
17/12 1.41666666667 0.0024531 False
41/29 1.41379310345 0.000420459 False
99/70 1.41428571429 7.21519e-5 False
239/169 1.41420118343 1.23789e-5 False
577/408 1.41421568627 2.1239e-6 False
1393/985 1.41421319797 3.64404e-7 True
3363/2378 1.41421362489 6.25218e-8 True

Uji: {q rasional : q^2 < 2} tidak punya maksimum
6625109/4684660 1.41421341143220639 True
149856914148961/105964844887080 1.41421349985133266 True
\end{verbatim}

Tiga hal terbaca dari keluaran itu. Galat Newton berhenti tercetak pada dua belas angka
di langkah 4, walaupun nilai sebenarnya masih turun sampai $6.4\times10^{-25}$: yang
membatasi lebar cetakan, bukan matematikanya. Ambang $10^{-6}$ dilewati tepat pada
$1393/985$, tempat kolom \texttt{True} mulai muncul. Kedua bilangan pada uji maksimum
juga \texttt{True}: keduanya rasional dan kuadratnya masih di bawah dua.

Kode itu tidak membuktikan apa pun. Ia memeriksa bahwa rumus di bab ini memberi angka
yang disebutkan dan bahwa arah pertidaksamaannya benar; buktinya tetap di lemma dan
teorema. Keluaran program adalah data, bukan argumen.

\begin{untukdosen}
Urutan mengajar yang saya sarankan: program pencarian akar di pembuka, lalu aksioma
lapangan dengan $0\cdot x=0$ dibuktikan di papan, lalu urutan, nilai mutlak, dan
kelengkapan. Jangan menunda kelengkapan; seluruh isi bab bergantung padanya.

Dua miskonsepsi yang sering muncul. Pertama, mahasiswa mengira supremum harus anggota
himpunan, sehingga mereka menolak $\sup\{1-1/n\}=1$. Kedua, mahasiswa mengira
kelengkapan bisa dibuktikan dari hukum lapangan dan urutan, padahal $\mathbb{Q}$
memenuhi keduanya dan gagal pada kelengkapan. Keduanya hilang setelah mahasiswa melihat
dua contoh itu berdampingan.

Aktivitas papan tulis yang berhasil: setiap mahasiswa mengambil $q=1.4$, menghitung
$h=(2-q^2)/(2q+2)$ dan memeriksa $(q+h)^2<2$; hasil yang berbeda-beda justru menunjukkan
rumusnya bekerja untuk setiap $q$. Tanda siap lanjut: mahasiswa dapat menjelaskan mengapa
$A$ tidak punya maksimum tanpa menghitung.
\end{untukdosen}

\section{Ringkasan}
\tingkat{dasar}

\begin{ringkasan}
\begin{itemize}
\item Bilangan real dilandasi tiga kelompok aksioma: lapangan, urutan, kelengkapan.
      Yang pertama memberi aljabar, yang kedua perbandingan, yang ketiga menutup lubang.
\item Dari lapangan diturunkan $0\cdot x=0$, $(-x)y=-(xy)$, dan ketiadaan pembagi nol.
\item Dari urutan diturunkan aturan tanda, termasuk kuadrat tak negatif dan $1>0$.
\item Nilai mutlak mengubah jarak menjadi pertidaksamaan: $|x|\le c$ setara dengan
      $-c\le x\le c$, dan $|x+y|\le|x|+|y|$.
\item Kelengkapan menjamin supremum setiap himpunan tak kosong yang terbatas di atas,
      lalu menurunkan Sifat Archimedes dan kerapatan $\mathbb{Q}$.
\item $\mathbb{Q}$ memenuhi lapangan dan urutan tetapi tidak lengkap: $\{q: q^2<2\}$
      tidak punya maksimum; di $\mathbb{R}$ supremumnya $\sqrt{2}$.
\end{itemize}
\end{ringkasan}

\section{Latihan}
\tingkat{menengah}

\begin{latihan}
\begin{enumerate}
  \item Buktikan bahwa dalam lapangan terurut, $a<b$ setara dengan $-b<-a$, dan
        simpulkan bahwa mengalikan pertidaksamaan dengan $-1$ membalik arahnya.
  \item Buktikan $a<\frac{a+b}{2}<b$ untuk $a<b$, lalu jelaskan mengapa hasil itu tidak
        memerlukan kelengkapan sedangkan $\sup\{q:q^2<2\}$ memerlukannya.
  \item Hitung $\sup$ dan $\inf$ himpunan $S=\{2-1/n : n\in\mathbb{N}\}$, lalu periksa
        kedua syarat definisi. Tentukan juga apakah $S$ punya maksimum.
  \item Nilai klaim ini: ``Karena setiap bilangan positif lebih besar dari $1/n$ untuk
        suatu $n$, bilangan real positif terkecil tidak ada.'' Benar, salah, atau belum
        cukup bukti? Sebutkan bagian yang didukung bukti.
  \item Dengan $h=(2-q^2)/(2q+2)$, hitung $q+h$ untuk $q=1$ dan $q=7/5$, lalu periksa
        bahwa keduanya masih di dalam $A=\{q\in\mathbb{Q}:q^2<2\}$. Bandingkan besar
        kenaikannya.
  \item Tunjukkan bahwa jika $|x|\le 2$ dan $|y|\le 2$ maka
        $|x^3-y^3|\le 12\,|x-y|$, dan sebutkan di mana ketaksamaan segitiga dipakai.
  \item Buktikan bahwa jika $S$ terbuka maka komplemennya tertutup, dan sebaliknya,
        dengan memakai definisi persekitaran.
  \item Hitung lima langkah iterasi Newton untuk $\sqrt{3}$ mulai dari $x_0=2$, dan
        laporkan galat relatif tiap langkah. Periksa apakah nisbah $e_{k+1}/e_k^2$
        mendekat ke $1/2$.
  \item Tunjukkan $\bigcap_{n=1}^{\infty}(-1/n,1/n)=\{0\}$ dan jelaskan mengapa itu tidak
        bertentangan dengan ketertutupan irisan berhingga.
  \item Buktikan bahwa di antara dua bilangan real selalu ada bilangan irasional.
\end{enumerate}

\noindent\textbf{Petunjuk.} Butir 4: klaim itu memuat dua pernyataan berbeda, satu
tentang setiap $\varepsilon$ dan satu tentang keberadaan minimum. Butir 8: gunakan
$e_{k+1}=e_k^2\cdot\frac{\sqrt3}{2x_k}$ dan $x_k\ge\sqrt3$. Butir 10: pakai kerapatan
$\mathbb{Q}$ pada $x/\sqrt2$ dan $y/\sqrt2$, lalu kalikan kembali.
\end{latihan}

\begin{untukpeneliti}
Yang benar-benar terbukti: aksioma kelengkapan tidak bergantung pada lapangan dan urutan,
sebab $\mathbb{Q}$ memenuhi keduanya dan gagal pada yang ketiga. Keberadaan lapangan
terurut lengkap dibuktikan Rudin lewat potongan Dedekind, dan di sana juga dinyatakan
tanpa bukti bahwa setiap dua lapangan terurut lengkap isomorfik; klaim itu tidak dipakai
di sini.

Sumber primer: Rudin, \emph{Principles of Mathematical Analysis}, Bab 1 (Definisi 1.12,
Proposisi 1.18, Teorema 1.19--1.21, Apendiks tentang potongan); Abbott,
\emph{Understanding Analysis}, Bab 1 (Aksioma Kelengkapan dan akibatnya).

Yang masih terbuka: isi efektif iterasi Newton. Tanpa batas awal eksplisit,
$e_{k+1}\le e_k^2/2$ belum memberi jumlah langkah; itu memerlukan Bab 4 dan, untuk
fungsi umum, teori Kantorovich. Batas keberlakuan bab: topologi sebatas interval dan
persekitaran, dan $\mathbb{R}$ diandaikan, bukan dibangun.
\end{untukpeneliti}

\peta{Bab 3 pada Rudin, \emph{Principles of Mathematical Analysis}: Bab 1 (Ordered Sets,
Fields, The Real Field, The Least-Upper-Bound Property, Sifat Archimedes, kerapatan
$\mathbb{Q}$, keberadaan akar ke-$n$, Apendiks potongan Dedekind). Pada Abbott,
\emph{Understanding Analysis}: Bab 1 (The Axiom of Completeness dan Consequences of
Completeness); himpunan terbuka dan titik akumulasi dari Bab 3, tidak dikembangkan di
sini.}
